\chapter{Полный product-a6 коэффициент циклической ветви}
\label{ch:v7-full-product-a6-cycle-coefficient-gate}

\section{Явные слабые интертвинеры}

Редуцированный граф дал ненулевой формальный коэффициент простого
шестикромочного цикла. Теперь раскроем слабые индексы. В единичной
калибровке общего Хиггса положим
\begin{equation}
 H=\begin{pmatrix}0\\1\end{pmatrix},\qquad
 \widetilde H=i\sigma_2\overline H
 =\begin{pmatrix}1\\0\end{pmatrix}.
 \label{eq:v7-full-a6-higgs-frame}
\end{equation}
Тогда старое up-ребро $Q_L\to u_R$ использует слабую ковекторную часть
$\widetilde H^\dagger$, а $L_L\to e_R$ --- $H^\dagger$. Два singlet-ребра
$L_L\leftrightarrow Y_R$ и цветное $Y_R\leftrightarrow Q_L$ сохраняют
слабый индекс тождественно. Поэтому слабый множитель исходного цикла равен
\begin{equation}
 \boxed{\widetilde H^\dagger H=0.}
 \label{eq:v7-full-a6-up-overlap-zero}
\end{equation}
Это то же ортогональное тождество, которое ранее запретило прямые
межрёберные слова второй степени. В полном произведении оно действует и
внутри шестишаговой композиции.

\section{Исходный цикл исчезает}

Пусть $x$ обозначает одну цветовую компоненту ребра $X_Lu_R$, а $y$ ---
согласованную компоненту $Q_LY_R$; четыре фоновых ребра имеют единичную
норму. Явная самосопряжённая блочная матрица даёт точный ограниченный след
\begin{equation}
 \Tr\Phi_u^6
 =50+24x^2+42y^2+12x^4+24y^4
 +6x^2y^2+2x^6+4y^6.
 \label{eq:v7-full-a6-up-polynomial}
\end{equation}
В нём отсутствует билинейный член:
\begin{equation}
 [xy]\Tr\Phi_u^6=0.
 \label{eq:v7-full-a6-up-mixed-zero}
\end{equation}
Следовательно, коэффициент $\kappa$ тяжёлого блока, использованный в
предыдущем determinant-контроле для пары $(X_Lu_R,Q_LY_R)$, равен нулю в
полном слабом product-следе. Формальная gauge-инвариантность замкнутого
слова не гарантирует его ненулевость.

\section{Исчерпывающая локальная конкуренция циклов}

Полный строгий граф содержит четырнадцать простых неориентированных
шестикромочных циклов. На фоне трёх старых рёбер и четырёх gauge-singlet
амплитуд ровно три из них содержат ровно два нулевых поля и потому входят в
тяжёлый гессиан:
\begin{equation}
\begin{array}{c|c|c}
\text{ветвь}&\text{два нулевых ребра}&\text{слабый множитель}\\ \hline
u&(Q_LY_R,X_Lu_R)&\widetilde H^\dagger H=0\\
d&(Q_LY_R,X_Ld_R)&H^\dagger H=1\\
W&(L_LX_R,Y_Le_R)&\langle\ell,y\rangle
\end{array}
 \label{eq:v7-full-a6-three-quadratic-cycles}
\end{equation}
Таким образом, полный оператор не просто уничтожает исходный цикл: он
переносит цветную ветвь с $u_R$ на $d_R$ и одновременно открывает
конкурирующую бесцветную пару слабых дублетов.

Для down-ветви явный след равен
\begin{equation}
 \Tr\Phi_d^6
 =50+24x^2+42y^2+12xy+12x^4+24y^4
 +6x^2y^2+2x^6+4y^6,
 \label{eq:v7-full-a6-down-polynomial}
\end{equation}
поэтому
\begin{equation}
 [xy]\Tr\Phi_d^6=12.
 \label{eq:v7-full-a6-down-mixed}
\end{equation}
Для параллельных компонент слабой пары получаем
\begin{equation}
 \Tr\Phi_W^6
 =48+30(x^2+y^2)+12xy+12(x^4+y^4)+2(x^6+y^6),
 \label{eq:v7-full-a6-weak-polynomial}
\end{equation}
тогда как для ортогональных компонент билинейный член исчезает. Это
подтверждает, что $12xy$ является интертвинерным скалярным произведением,
а не артефактом простой смежности.

\section{Gaussian-знак и гессиан шестого следа}

На постоянном плоском фоне чистый Gaussian даёт потенциальный множитель
\begin{equation}
 -\frac16\Tr\Phi^6.
 \label{eq:v7-full-a6-gaussian-factor}
\end{equation}
Следовательно, билинейные вклады трёх ветвей равны соответственно
\begin{equation}
 0,\qquad -2xy,\qquad -2xy.
 \label{eq:v7-full-a6-gaussian-mixed-terms}
\end{equation}
При сохранении ориентации это знак центрального минимума $W_C=I$, а не
$-I$. Однако он относится уже к down-циклу.

Полные ограниченные гессианы $\Tr\Phi^6$ в начале тяжёлых пар равны
\begin{align}
 H_u&=\begin{pmatrix}48&0\\0&84\end{pmatrix},\nonumber\\
 H_d&=\begin{pmatrix}48&12\\12&84\end{pmatrix},\qquad
 \Spec H_d=\{66-6\sqrt{13},66+6\sqrt{13}\},\nonumber\\
 H_W&=\begin{pmatrix}60&12\\12&60\end{pmatrix},\qquad
 \Spec H_W=\{48,72\}.
 \label{eq:v7-full-a6-heavy-hessians}
\end{align}
Gaussian-$a_6$ умножает их на отрицательное число. Поэтому устойчивость
нельзя решать по одному смешанному моному: положительные $a_2/a_4$ части и
точный экспоненциальный остаток обязательны.

\begin{theorem}[Обнуление исходного цикла и перенос ветви]
После раскрытия стандартных слабых интертвинеров коэффициент исходного
цикла через $u_R$ тождественно равен нулю из-за
$\widetilde H^\dagger H=0$. Из четырнадцати простых шестикромочных циклов
ровно три входят квадратично по тяжёлым полям около singlet-вакуума. Один
обнулён, down-цикл через $X_Ld_R$ имеет коэффициент двенадцать, и ещё один
коэффициент двенадцать связывает пару слабых дублетов. Следовательно,
прежняя однопарная виртуальная модель не является редукцией полного
product-$a_6$.
\end{theorem}

\begin{nogocriterion}[Запрет редуцированного циклического коэффициента]
Нельзя переносить коэффициент $12$ скалярной смежности на полный цикл через
$u_R$ без слабого следа. Нельзя также молча заменить $X_Lu_R$ на прежде
конкурирующее $X_Ld_R$: полный оператор одновременно открывает слабую пару
$(L_LX_R,Y_Le_R)$, и обе выжившие ветви должны пройти один общий точный
Gaussian-гессиан.
\end{nogocriterion}

\begin{protocol}
\begin{enumerate}
 \item простых шестикромочных циклов полного графа: \textbf{$14$};
 \item циклов, квадратичных по нулевым полям около singlet-вакуума:
 \textbf{$3$};
 \item слабый множитель исходного up-цикла:
 \textbf{$\widetilde H^\dagger H=0$};
 \item полный коэффициент исходного up-цикла: \textbf{$0$};
 \item коэффициент down-цикла: \textbf{$12$};
 \item коэффициент слабой конкурирующей пары: \textbf{$12$};
 \item Gaussian-билинейный коэффициент выживших ветвей: \textbf{$-2$};
 \item прежний determinant-параметр $\kappa$ для up-пары: \textbf{$0$};
 \item полный product-$a_6$ проход исходного цикла: \textbf{нет};
 \item down-ветвь: \textbf{открыта};
 \item слабая конкуренция: \textbf{обязательна};
 \item итог: \textbf{строгий no-go исходного цикла и перенос ветви};
 \item следующий гейт: \textbf{конкуренция weak-aligned down-цикла}.
\end{enumerate}
\end{protocol}

Машинный сертификат сохранён в
\path{s2t_v7_full_product_a6_cycle_coefficient_gate_results.json}.

\begin{thebibliography}{9}
\bibitem{v7-full-a6-vassilevich}
D.~V.~Vassilevich, \emph{Heat Kernel Expansion: User's Manual},
Physics Reports 388 (2003), 279--360; arXiv:hep-th/0306138.
\bibitem{v7-full-a6-dimension-six}
A.~Devastato, F.~Lizzi, C.~Valcarcel Flores, D.~Vassilevich,
\emph{Unification of Coupling Constants, Dimension Six Operators and the
Spectral Action}, International Journal of Modern Physics A 30 (2015),
1550033; arXiv:1410.6624.
\bibitem{v7-full-a6-quiver}
C.~I.~P\'erez-S\'anchez,
\emph{The Spectral Action on Quivers}, arXiv:2401.03705.
\bibitem{v7-full-a6-ncg-sm}
A.~Devastato, \emph{Spectral Noncommutative Geometry, Standard Model and
all that}, arXiv:1906.09583.
\end{thebibliography}